\section{Additional Empirical Evidence}
\label{sec:app-empirical}

\subsection{Sample composition and specification history}
\label{sec:app-sample}

% AUTO-GENERATED by pipeline render-paper5-descriptives -- do not edit by hand.

\begin{table}[H]

\centering

\small
\begin{tabularx}{\linewidth}{Xrrrrr}
\toprule
Sector & $n$ & Articles & Volatility & $W_2$ within & $W_2$ cross \\
\midrule
Communication Services & 11 & 128.0 & 0.0239 & 1.042 & 1.060 \\
Consumer Cyclical & 18 & 128.0 & 0.0281 & 1.033 & 1.053 \\
Consumer Defensive & 7 & 128.0 & 0.0176 & 1.010 & 1.058 \\
Energy & 1 & 128.0 & 0.0386 & \textemdash & 1.045 \\
Financial Services & 1 & 128.0 & 0.0272 & \textemdash & 1.070 \\
Healthcare & 14 & 128.0 & 0.0233 & 1.025 & 1.065 \\
Industrials & 7 & 128.0 & 0.0240 & 1.030 & 1.057 \\
Technology & 38 & 128.0 & 0.0265 & 0.998 & 1.050 \\
Utilities & 3 & 128.0 & 0.0159 & 0.882 & 1.058 \\
\midrule
All sectors & 100 & 128.0 & 0.0251 & 1.008 & 1.056 \\
\bottomrule
\end{tabularx}

\caption{Per-sector composition of the news frame and its balanced-cloud sample: firm count, mean article count per firm, mean daily return volatility, and mean within- versus cross-sector rooted $W_2$ distance under the frozen Qwen3-Embedding-8B geometry. Firm counts, volatility, and distance summaries use the same estimation sample. The Articles column reports the constant balanced-cloud size after per-firm truncation, not raw frame coverage, and therefore carries no information about a firm's true news volume. Volatility is daily and unannualized. A dash marks a statistic undefined for that sector, such as within-sector distance for a singleton.}

\label{tab:p5-sample-descriptives}

\end{table}


\Cref{tab:p5-sample-descriptives} records the news frame's industry
composition and its return and text dispersion.
Sector sizes are heavily unbalanced, with Technology holding roughly a third of
the frame and two sectors holding fewer than three firms, so peer sets built
from the text geometry inherit that imbalance.
Mean within-sector $W_2$ distance is below the corresponding cross-sector
distance in every sector holding more than one firm: the frozen geometry
contains industry structure before any return is used.

The reported specification expands the priced universe from an earlier 52-firm
specification to \ppvalue{n-tickers-priced} firms and ends the return window
on 6 April 2026, the final common date before HOLX's take-private.
\Cref{tab:p5-universe-horizon-bridge} bridges the three observed
specifications.
At 52 firms, extending the common-date window from 6 April to 15 July 2026
moves $\hat\rho$ from
$\ppnum[3]{bridge-reference-52-short-rho}$ to
$\ppnum[3]{bridge-reference-52-long-rho}$, a difference of
$\ppnum[3]{bridge-horizon-rho-difference}$.
At the short horizon, expanding the universe from 52 to 100 firms moves
$\hat\rho$ from $\ppnum[3]{bridge-reference-52-short-rho}$ to
$\ppnum[3]{bridge-current-100-short-rho}$, a difference of
$\ppnum[3]{bridge-universe-rho-difference}$.
Both are arithmetic differences reported without contrast-level inference.
The universe comparison changes the frozen operator as well as the roster, and
its complete-case intersection has one fewer date.
The 52 shared return histories are identical through the common cutoff, so
data revision within those histories does not explain the comparison.

% AUTO-GENERATED -- do not edit by hand.
\begin{table}[H]
\centering
\small
\begin{tabular}{lrrrl}
\toprule
Specification & Firms & End date & Trading days & $\hat\rho$ [95\% CI] \\
\midrule
Historical reference & 52 & 15 Jul 2026 & 885 & 0.776 [0.743, 0.807] \\
Horizon-matched & 52 & 6 Apr 2026 & 816 & 0.771 [0.735, 0.806] \\
Current universe & 100 & 6 Apr 2026 & 815 & 0.843 [0.815, 0.869] \\
\bottomrule
\end{tabular}
\caption{Sequential universe-and-horizon bridge for the barycentric $W^\flat$ single-field SAR QMLE in the raw-return, common-intercept specification (\ppAppendixRef{sec:app-unadjusted}); it is not the primary factor-conditioned estimate. Every row starts on 3 January 2023. The historical row uses the recovered 52-firm operator and the long return window. The horizon-matched row holds that operator and universe fixed while shortening the return endpoint; its interval is newly estimated with the joint-date stationary bootstrap. The current-universe row expands to the 100-firm operator at the same endpoint, so it changes both the roster and the operator; the first and third intervals come from the reported source results. Complete-case trading-day counts can differ across universes.}
\label{tab:p5-universe-horizon-bridge}
\end{table}


\subsection{The complete barycentric field}
\label{sec:app-field-anatomy}

\Cref{fig:wflat-barycentric-mixtures} displays the complete field.
Each row is a target firm and each column a candidate source; color marks the
five largest actual coefficients in each row without renormalizing them.
The upper panel reports each source's total incoming mass across target rows,
and the right-hand panels report each target's effective source count and the
share of its row captured by the five displayed cells.

\begin{figure}[htbp]
  \centering
  \ppfigure{1}{w2_barycentric_mixture_weights}
  \caption{Illustrative barycentric interaction field from the target-anchored
    Wasserstein barycentric reconstruction. Each row is the actual
    leave-one-out solution $W^\flat_{i\cdot}$ to \eqref{eq:w2-barycentre} over
    the frozen 2018--2022 article clouds: weights are nonnegative, sum to one,
    and assign zero self-weight. The heat map shows each target's five largest
    \emph{actual} coefficients without renormalizing them; the remaining
    positive coefficients still belong to $W^\flat$ but are not colored. The
  upper and right-hand diagnostics are calculated from the complete field.}
  \label{fig:wflat-barycentric-mixtures}
  \figalttext{Rows index target firms and columns index source firms. Colored
    cells mark each target's five largest actual barycentric weights; their
    scattered, row-specific locations show that the field is directed and
    heterogeneous. Marginal bar charts show heterogeneous incoming mass,
    effective source counts, and the fraction of each full row contained in the
  displayed top five. Uncolored positive weights remain in the operator.}
\end{figure}

Both concentration within rows and total incoming mass across rows vary across
firms, and the colored locations are scattered and row specific: the field is
directed and heterogeneous rather than a block-diagonal industry map.
\Cref{tab:p5-peer-interpretation} lists actual reconstruction weights $W^\flat_{ij}$,
pairwise-distance ranks, sector and industry matches, and persistent
co-mention status for AMD and for the targets with minimum, median, and maximum
effective source count.
A large weight records a peer's usefulness within the target's reconstructing
combination, which need not coincide with its pairwise rank.

% AUTO-GENERATED -- do not edit by hand.
\begin{table}[H]
\centering
\scriptsize
\setlength{\tabcolsep}{2.5pt}
\begin{tabular}{ll l r r c c c}
\toprule
Target & Selection role & Peer & $W^\flat_{ij}$ & W2 rank & Same sector & Same industry & Co-mention \\
\midrule
AMD & declared AMD & NVDA & 0.1225 & 1 & Yes & Yes & Yes \\
AMD & declared AMD & INTC & 0.0987 & 4 & Yes & Yes & Yes \\
AMD & declared AMD & MU & 0.0878 & 3 & Yes & Yes & Yes \\
AMD & declared AMD & AMAT & 0.0840 & 2 & Yes & No & Yes \\
AMD & declared AMD & ASML & 0.0688 & 6 & Yes & No & Yes \\
\addlinespace[1.5pt]
HON & maximum effective source count & CSCO & 0.0534 & 1 & No & No & Yes \\
HON & maximum effective source count & UAL & 0.0502 & 51 & Yes & No & Yes \\
HON & maximum effective source count & KHC & 0.0427 & 24 & No & No & No \\
HON & maximum effective source count & PEP & 0.0405 & 7 & No & No & Yes \\
HON & maximum effective source count & AAL & 0.0387 & 60 & Yes & No & Yes \\
\addlinespace[1.5pt]
JBHT & median effective source count & CSX & 0.1473 & 1 & Yes & No & Yes \\
JBHT & median effective source count & PCAR & 0.0629 & 3 & Yes & No & No \\
JBHT & median effective source count & UAL & 0.0524 & 43 & Yes & No & Yes \\
JBHT & median effective source count & CTSH & 0.0458 & 2 & No & No & No \\
JBHT & median effective source count & AEP & 0.0387 & 5 & No & No & No \\
\addlinespace[1.5pt]
XEL & minimum effective source count & EXC & 0.2443 & 1 & Yes & Yes & Yes \\
XEL & minimum effective source count & AEP & 0.2173 & 2 & Yes & Yes & Yes \\
XEL & minimum effective source count & ENPH & 0.0797 & 3 & No & No & Yes \\
XEL & minimum effective source count & FANG & 0.0541 & 5 & No & No & No \\
XEL & minimum effective source count & CSX & 0.0409 & 4 & No & No & No \\
\bottomrule
\end{tabular}
\caption{Descriptive W2 peer interpretation panel. Targets are AMD and deterministic minimum, median, and maximum effective source-count firms, with lexical tie-breaking. Each row lists one of the target's five largest actual reconstruction weights; the distance rank is the target's pairwise W2 rank. Sector, industry, and co-mention columns report metadata associated with each candidate.}
\label{tab:p5-peer-interpretation}
\end{table}


\subsection{The common-factor ladder and its information criteria}
\label{sec:app-information-criteria}

\Cref{tab:factor-set-sensitivity} repeats the primary specification under
each common-factor control set on the same dates.
The spatial coefficient is stable across the ladder and every bootstrap
interval excludes zero; adding momentum lowers $\hat\rho$ by about one
percentage point within either factor family.

\begin{table}[H]
  \centering
  \caption{Common-factor sensitivity for the barycentric field $W^\flat$.
    Each row profiles firm-specific intercepts and loadings on the indicated
    factors. The empirical 95\% intervals use 2,000 paired joint-date
    stationary-bootstrap refits. The held-out score fits all parameters on
    2018--2022 and evaluates the mean Gaussian quasi-log score per firm-day on
    2023--2026 returns; higher is better. FF5 plus momentum is the prespecified
  primary model.}
  \label{tab:factor-set-sensitivity}
  \begin{tabular}{lrrrr}
    \toprule
    Factor controls & $K$ & $\hat\rho$ & 95\% interval & Held-out score \\
    \midrule
    FF3 &
    \ppvalue{transport-factor-ff3-n-factors} &
    \ppnum[3]{transport-factor-ff3-w2-rho} &
    $[\ppnum[3]{transport-factor-ff3-w2-rho-ci-lower},
    \ppnum[3]{transport-factor-ff3-w2-rho-ci-upper}]$ &
    \ppnum[6]{transport-factor-ff3-w2-predictive-score} \\
    FF3 plus momentum &
    \ppvalue{transport-factor-ff3-mom-n-factors} &
    \ppnum[3]{transport-factor-ff3-mom-w2-rho} &
    $[\ppnum[3]{transport-factor-ff3-mom-w2-rho-ci-lower},
    \ppnum[3]{transport-factor-ff3-mom-w2-rho-ci-upper}]$ &
    \ppnum[6]{transport-factor-ff3-mom-w2-predictive-score} \\
    FF5 &
    \ppvalue{transport-factor-ff5-n-factors} &
    \ppnum[3]{transport-factor-ff5-w2-rho} &
    $[\ppnum[3]{transport-factor-ff5-w2-rho-ci-lower},
    \ppnum[3]{transport-factor-ff5-w2-rho-ci-upper}]$ &
    \ppnum[6]{transport-factor-ff5-w2-predictive-score} \\
    FF5 plus momentum (primary) &
    \ppvalue{transport-factor-ff5-mom-n-factors} &
    \ppnum[3]{transport-factor-ff5-mom-w2-rho} &
    $[\ppnum[3]{transport-factor-ff5-mom-w2-rho-ci-lower},
    \ppnum[3]{transport-factor-ff5-mom-w2-rho-ci-upper}]$ &
    \ppnum[6]{transport-factor-ff5-mom-w2-predictive-score} \\
    \bottomrule
  \end{tabular}
\end{table}

The held-out score is highest under FF3, and working information criteria rank
the ladder differently according to their penalty; FF5 plus momentum remains
primary because it conditions on the broadest prespecified set of observed
factors.

\begin{table}[H]
  \centering
  \caption{Working information criteria for the barycentric factor ladder.
    The first panel reports the maximized post-corpus quasi-log-likelihood and
    AIC. The second reports HQC and BIC using either multivariate dates, $T$,
    or firm-days, $NT$, as the effective sample size. Lower information
  criteria are better.}
  \label{tab:factor-information-criteria}
  \begin{tabular}{lrrr}
    \toprule
    Factor controls & Parameters & QLL & AIC \\
    \midrule
    FF3 &
    \ppvalue{spatial-value-factor-ff3-parameter-count} &
    \ppnum[1]{spatial-value-factor-ff3-loglik} &
    \ppnum[1]{spatial-value-factor-ff3-aic} \\
    FF3 plus momentum &
    \ppvalue{spatial-value-factor-ff3-mom-parameter-count} &
    \ppnum[1]{spatial-value-factor-ff3-mom-loglik} &
    \ppnum[1]{spatial-value-factor-ff3-mom-aic} \\
    FF5 &
    \ppvalue{spatial-value-factor-ff5-parameter-count} &
    \ppnum[1]{spatial-value-factor-ff5-loglik} &
    \ppnum[1]{spatial-value-factor-ff5-aic} \\
    FF5 plus momentum &
    \ppvalue{spatial-value-factor-ff5-mom-parameter-count} &
    \ppnum[1]{spatial-value-factor-ff5-mom-loglik} &
    \ppnum[1]{spatial-value-factor-ff5-mom-aic} \\
    \bottomrule
  \end{tabular}

  \begin{tabular}{lrrrr}
    \toprule
    Factor controls & $\mathrm{HQC}_T$ & $\mathrm{HQC}_{NT}$ &
    $\mathrm{BIC}_T$ & $\mathrm{BIC}_{NT}$ \\
    \midrule
    FF3 &
    \ppnum[1]{spatial-value-factor-ff3-hqc-dates} &
    \ppnum[1]{spatial-value-factor-ff3-hqc-firm-days} &
    \ppnum[1]{spatial-value-factor-ff3-bic-dates} &
    \ppnum[1]{spatial-value-factor-ff3-bic-firm-days} \\
    FF3 plus momentum &
    \ppnum[1]{spatial-value-factor-ff3-mom-hqc-dates} &
    \ppnum[1]{spatial-value-factor-ff3-mom-hqc-firm-days} &
    \ppnum[1]{spatial-value-factor-ff3-mom-bic-dates} &
    \ppnum[1]{spatial-value-factor-ff3-mom-bic-firm-days} \\
    FF5 &
    \ppnum[1]{spatial-value-factor-ff5-hqc-dates} &
    \ppnum[1]{spatial-value-factor-ff5-hqc-firm-days} &
    \ppnum[1]{spatial-value-factor-ff5-bic-dates} &
    \ppnum[1]{spatial-value-factor-ff5-bic-firm-days} \\
    FF5 plus momentum &
    \ppnum[1]{spatial-value-factor-ff5-mom-hqc-dates} &
    \ppnum[1]{spatial-value-factor-ff5-mom-hqc-firm-days} &
    \ppnum[1]{spatial-value-factor-ff5-mom-bic-dates} &
    \ppnum[1]{spatial-value-factor-ff5-mom-bic-firm-days} \\
    \bottomrule
  \end{tabular}
\end{table}

Because each added factor introduces one loading for every firm, the working
information criteria impose materially different penalties from the held-out
score comparison.
Working AIC and both HQC conventions are lowest for FF5 plus momentum, whereas
both BIC conventions are lowest for FF3 plus momentum; the ranking depends on
the penalty and the effective sample size.
The criteria are
$\mathrm{AIC}=-2\ell+2k$,
$\mathrm{HQC}_T=-2\ell+2k\log\log T$,
$\mathrm{HQC}_{NT}=-2\ell+2k\log\log(NT)$,
$\mathrm{BIC}_T=-2\ell+k\log T$, and
$\mathrm{BIC}_{NT}=-2\ell+k\log(NT)$, using dates and firm-days,
respectively, as the effective sample size; lower values are preferred.
For $K$ factors, $k=N(K+1)+1$ under the spherical factor-only model and
$N(K+1)+2$ under a spatial alternative; the diagonal-variance models have
$N(K+1)+N$ and $N(K+1)+N+1$ parameters.
Momentum additions are nested within each factor family, whereas the official
FF3 and FF5 files differ in their SMB construction.

\subsection{Firm-specific innovation variances}
\label{sec:app-diagonal}

\begin{table}[H]
  \centering
  \caption{Spatial-value robustness with firm-specific innovation variances.
    QLL gains, restricted-null QLR tests, and held-out score gains compare each
  field with the matched diagonal-variance factor-only model.}
  \label{tab:spatial-value-diagonal}
  \begin{tabular}{lrrrr}
    \toprule
    Field & $\hat\rho$ & QLL gain & QLR $p$ & Held-out gain \\
    \midrule
    Target-anchored W1 &
    \ppnum[3]{spatial-value-diagonal-w1-rho} &
    \ppnum[1]{spatial-value-diagonal-w1-versus-none-loglik-gain} &
    \ppnum[4]{spatial-value-diagonal-w1-versus-none-qlr-pvalue} &
    $\ppnum[4]{spatial-value-diagonal-w1-versus-none-held-out-score-gain}$
    $[\ppnum[4]{spatial-value-diagonal-w1-versus-none-held-out-score-gain-ci-lower},
    \ppnum[4]{spatial-value-diagonal-w1-versus-none-held-out-score-gain-ci-upper}]$
    \\
    Target-anchored W2 &
    \ppnum[3]{spatial-value-diagonal-w2-rho} &
    \ppnum[1]{spatial-value-diagonal-w2-versus-none-loglik-gain} &
    \ppnum[4]{spatial-value-diagonal-w2-versus-none-qlr-pvalue} &
    $\ppnum[4]{spatial-value-diagonal-w2-versus-none-held-out-score-gain}$
    $[\ppnum[4]{spatial-value-diagonal-w2-versus-none-held-out-score-gain-ci-lower},
    \ppnum[4]{spatial-value-diagonal-w2-versus-none-held-out-score-gain-ci-upper}]$
    \\
    \bottomrule
  \end{tabular}
\end{table}

With firm-specific variances, each spatial model also has lower working AIC,
HQC, and BIC than its matched factor-only benchmark, and the reduced-form
coefficient identity holds to within its $10^{-10}$ absolute and relative
tolerance.

\subsection{Covariance-sensitive portfolio diagnostic}
\label{sec:app-gmv}

Proposition~\ref{prop:factor-beta-invariance} implies that the economic
difference across fields lies in the covariance assigned to factor-residual
returns. For each fitted field, the stored total covariance is
\begin{equation}
  \widehat{\Sigma}_{r,m}
  =\widehat{B}\widehat{\Sigma}_{f}\widehat{B}^{\prime}
  + A_m^{-1}\widehat{\Omega}_m A_m^{-\prime},
  \qquad A_m=I-\hat\rho_m W_m,
  \label{eq:app-gmv-covariance}
\end{equation}
where $m$ indexes the factor-only, W1, and W2 specifications and
$\widehat{\Omega}_m$ is the constant-correlation Ledoit--Wolf estimate of the
factor-residualized spatial innovations. The factor-only row sets
$\hat\rho_m=0$ and uses the same residual estimator. Thus the covariance inputs
combine common-factor covariance with the same estimator-defined local
factor-residual covariance rule across candidates.
We translate these covariance estimates into a fixed-weight allocation
diagnostic. The weights are formed on the factor-complete 2018--2022 panel,
frozen at the formation vintage, and evaluated on 2023--2026 excess returns
under a fully invested long-only constraint with a 15\% single-name cap.

% AUTO-GENERATED by pipeline render-paper5-gmv-validation-table -- do not edit by hand.
\begin{table}[H]
  \centering
  \scriptsize
  \setlength{\tabcolsep}{3.0pt}
  \begin{tabular}{lrrrrr}
    \toprule
    Covariance input & Realized variance & Ann. vol. & Rel. to factor only & Eff. $N$ & Max. wt. \\
    \midrule
Factor only + LW residual covariance & 0.5441 & 11.71\% & 100.0\% & 9.76 & 15.0\% \\
Target-anchored W1 & 0.5459 & 11.73\% & 100.3\% & 9.82 & 15.0\% \\
Target-anchored W2 & 0.5455 & 11.72\% & 100.3\% & 9.83 & 15.0\% \\
Sample GMV & 0.5428 & 11.70\% & 99.8\% & 9.57 & 15.0\% \\
Direct Ledoit--Wolf GMV & 0.5395 & 11.66\% & 99.2\% & 9.96 & 15.0\% \\
Equal weight & 1.1606 & 17.10\% & 213.3\% & 100.00 & 1.0\% \\
    \bottomrule
  \end{tabular}
  \caption{Covariance-sensitive portfolio diagnostic. Candidate weights are formed from factor-complete 2018--2022 excess returns and frozen at the 2022 formation vintage. They are long-only, fully invested, and subject to a 15\% single-name cap; realized risk is evaluated on constant-weight excess-return mixtures over 2023--2026. The factor-only, W1, and W2 rows use diagonal innovation variances with common constant-correlation residual shrinkage. Realized variance is scaled by $10^4$; lower values are better.}
  \label{tab:p5-gmv-validation}
\end{table}


Under this fixed-weight GMV convention, W1 and W2 do not reduce realized
variance relative to the factor-only covariance benchmark; direct
Ledoit--Wolf shrinkage is slightly lower. The exhibit is therefore a diagnostic
of the residual-covariance interpretation, while the paper's primary evidence
remains the held-out joint-density comparison.

\subsection{Transport order at the article level and jointly}
\label{sec:app-transport-order}

The transport-order equivalence holds across the factor ladder.
The W2-minus-W1 coefficient difference ranges from
\ppnum[6]{transport-factor-ff3-mom-delta-rho} to
\ppnum[6]{transport-factor-ff5-delta-rho}, and the held-out score difference
from \ppnum[6]{transport-factor-ff5-delta-predictive-score} to
\ppnum[6]{transport-factor-ff3-delta-predictive-score} per firm-day.
After projection on the intercept and factors, the induced peer-return signals
correlate at
\ppnum[6]{transport-factor-residualized-signal-correlation} with a normalized
difference of
\ppnum[3]{transport-factor-residualized-relative-signal-difference-sd}, so the
factor-residualized quantities leave the weak-separation classification
unchanged.

At full dimension the two transport orders send
\ppnum[0]{transport-switch-d4096-switched-atoms} target atoms to different peer
atoms.
Scoring both assignments on the same chord distance, the mean W2-minus-W1
difference is \ppnum[4]{transport-switch-d4096-max-difference} in the longest
matched distance and \ppnum[4]{transport-switch-d4096-p95-difference} at its
95th percentile, against
$+\ppnum[4]{transport-switch-d4096-switched-difference}$ in mean matched
distance on the reassigned atoms.
Quadratic transport thus buys a shorter displacement tail with slightly longer
typical matches, which is the economic content of squaring the ground cost.
W2 gives a strictly shorter longest match in
\ppnum[3]{transport-switch-d4096-max-reduction-share} of firm pairs, at a mean
difference of \ppnum[3]{transport-switch-d4096-max-reduction-given-lower}, and
a strictly longer one in
\ppnum[3]{transport-switch-d4096-max-increase-share}, at
$+\ppnum[3]{transport-switch-d4096-max-increase-given-higher}$; the remaining
pairs share the same longest match.
The pattern strengthens monotonically down the dimension ladder, reaching
\ppnum[4]{transport-switch-d64-max-difference} in the longest match and a
shorter-tail share of \ppnum[3]{transport-switch-d64-max-reduction-share} at 64
dimensions.

After observing the standalone comparison, we estimated both fields jointly,
entering each supplied matrix unchanged:

\begin{equation}\label{eq:factor-two-field-transport}
  (I-\rho^\flat_1 W^\flat_1-\rho^\flat W^\flat)r_t^e
  =\alpha+A f_t+\varepsilon_t,
  \qquad
  \rho^\flat_1,\rho^\flat\geq 0,
  \quad \rho^\flat_1+\rho^\flat<1.
\end{equation}

At every candidate $(\rho^\flat_1,\rho^\flat)$ the specification re-profiles the same
firm intercepts, factor loadings, and innovation variance as
\eqref{eq:factor-spatial-lag}, and the restrictions $\rho^\flat=0$ and $\rho^\flat_1=0$
recover the single-field models.
The incremental statistic for W2 given W1 is
\[
  \operatorname{QLR}_{2\mid 1}
  =2\left\{
    \ell^Q(\hat\rho^\flat_1,\hat\rho^\flat)
    -\ell^Q(\hat\rho^{\flat,0}_1,0)
  \right\},
\]
with the reverse statistic defined analogously.
Because each null places a coefficient on the boundary, reference distributions
come from 2,000 restricted-null stationary-block innovation resamples generated
through the restricted spatial inverse, with the factor nuisance space projected
anew before both models are refit.

\begin{table}[H]
  \centering
  \caption{Factor-conditioned joint W1 and W2 specification. Both
    supplied fields
    enter unchanged under nonnegative coefficients whose sum is below one.
    Intervals are descriptive 95\% bootstrap quantiles, and the
    coefficient-zero share is the fraction of draws in which that coefficient is
    estimated at zero; both use 2,000 joint-date stationary-bootstrap refits.
  The specification was added after the standalone transport-order comparison.}
  \label{tab:transport-order-joint}
  \begin{tabular}{lrrr}
    \toprule
    Quantity & Estimate & 95\% interval & Coefficient-zero share \\
    \midrule
    W1 coefficient $\rho^\flat_1$ &
    \ppnum[3]{transport-joint-w1-rho} &
    $[\ppnum[3]{transport-joint-rho-w1-ci-lower},
    \ppnum[3]{transport-joint-rho-w1-ci-upper}]$ &
    \ppnum[3]{transport-joint-w2-only-boundary-share} \\
    W2 coefficient $\rho^\flat$ &
    \ppnum[3]{transport-joint-w2-rho} &
    $[\ppnum[3]{transport-joint-rho-w2-ci-lower},
    \ppnum[3]{transport-joint-rho-w2-ci-upper}]$ &
    \ppnum[3]{transport-joint-w1-only-boundary-share} \\
    Total feedback $\rho^\flat_1+\rho^\flat$ &
    \ppnum[3]{transport-joint-rho-total} &
    $[\ppnum[3]{transport-joint-rho-total-ci-lower},
    \ppnum[3]{transport-joint-rho-total-ci-upper}]$ & --- \\
    \bottomrule
  \end{tabular}
\end{table}

The joint estimates are $\hat\rho^\flat_1=
\ppnum[3]{transport-joint-w1-rho}$ and $\hat\rho^\flat=
\ppnum[3]{transport-joint-w2-rho}$, so allowing W2 adds
\ppnum[2]{transport-joint-loglik-gain-w1} quasi-log-likelihood units beyond
W1 alone.
The bootstrap places \ppnum[3]{transport-joint-w1-only-boundary-share} of draws
on the W1-only boundary, where $\rho^\flat=0$, and the two channel estimates
correlate at
\ppnum[3]{transport-joint-channel-correlation}: total feedback is precisely
positive, while its allocation between two nearly collinear fields is weakly
determined.
For W2 conditional on W1,
$\operatorname{QLR}_{2\mid 1}=
\ppnum[3]{transport-joint-qlr-w2-given-w1-statistic}$ with working-null
$p=\ppnum[4]{transport-joint-qlr-w2-given-w1-pvalue}$; for W1 conditional on
W2, $\operatorname{QLR}_{1\mid 2}=
\ppnum[3]{transport-joint-qlr-w1-given-w2-statistic}$ with
$p=\ppnum[4]{transport-joint-qlr-w1-given-w2-pvalue}$.
The in-sample nested fit therefore favors W1, while the held-out scores remain
predictively equivalent.

\subsection{Raw-return, common-intercept estimates}
\label{sec:app-unadjusted}

\begin{table}[H]
  \centering
  \caption{Pooled 2023--2026 single-field estimates in the
    raw-return, common-intercept
    specification. The models use
    raw returns and a common intercept. The 95\% intervals use 2,000
    joint-date stationary-bootstrap refits, and
    $\hat\rho=\hat\lambda/(1+\hat\lambda)$. The final column reports the
  maximized conditional quasi-log-likelihood on the common return sample.}
  \label{tab:qmle-pooled}
  \begin{tabular}{lrrrr}
    \toprule
    Interaction field & $\hat\lambda$ & 95\% interval & Implied $\hat\rho$
    & $\ell^Q$ \\
    \midrule
    Barycentric field ($W^\flat$) &
    \ppnum[2]{post-pooled-2023-2026-w-flat-lambda} &
    $[\ppnum[2]{post-pooled-2023-2026-w-flat-lambda-ci-lower},
    \ppnum[2]{post-pooled-2023-2026-w-flat-lambda-ci-upper}]$ &
    \ppnum[3]{post-pooled-2023-2026-w-flat-rho} &
    \ppnum[1]{post-pooled-2023-2026-w-flat-loglik} \\
    RBF--Wasserstein diffusion &
    \ppnum[2]{post-pooled-2023-2026-w-h-lambda} &
    $[\ppnum[2]{post-pooled-2023-2026-w-h-lambda-ci-lower},
    \ppnum[2]{post-pooled-2023-2026-w-h-lambda-ci-upper}]$ &
    \ppnum[3]{post-pooled-2023-2026-w-h-rho} &
    \ppnum[1]{post-pooled-2023-2026-w-h-loglik} \\
    Persistent news co-mentions &
    \ppnum[2]{post-pooled-2023-2026-w-co-mentions-lambda} &
    $[\ppnum[2]{post-pooled-2023-2026-w-co-mentions-lambda-ci-lower},
    \ppnum[2]{post-pooled-2023-2026-w-co-mentions-lambda-ci-upper}]$ &
    \ppnum[3]{post-pooled-2023-2026-w-co-mentions-rho} &
    \ppnum[1]{post-pooled-2023-2026-w-co-mentions-loglik} \\
    \midrule
    Equal active support &
    \ppnum[2]{post-pooled-2023-2026-equal-support-lambda} &
    $[\ppnum[2]{post-pooled-2023-2026-equal-support-lambda-ci-lower},
    \ppnum[2]{post-pooled-2023-2026-equal-support-lambda-ci-upper}]$ &
    \ppnum[3]{post-pooled-2023-2026-equal-support-rho} &
    \ppnum[1]{post-pooled-2023-2026-equal-support-loglik} \\
    \bottomrule
  \end{tabular}
\end{table}

In the raw-return, common-intercept specification, the barycentric field has
pooled adjustment index
$\hat\lambda=\ppnum[2]{post-pooled-2023-2026-w-flat-lambda}$, with interval
$[\ppnum[2]{post-pooled-2023-2026-w-flat-lambda-ci-lower},
\ppnum[2]{post-pooled-2023-2026-w-flat-lambda-ci-upper}]$, and feedback
coefficient $\hat\rho=\ppnum[3]{post-pooled-2023-2026-w-flat-rho}$.
When the identification condition in \Cref{sec:estimand} holds, this index
corresponds to an adjustment criterion that places about
\ppnum[2]{post-pooled-2023-2026-w-flat-lambda} times as much weight on
misalignment with the barycentric peer average as on departure from
stand-alone exposure; the interval determines the sign more sharply than the
magnitude.
Equal weighting of the same active support and pairwise RBF proximity both
yield positive feedback and lower quasi-log-likelihood than the barycentric
field, and the news-link field attains higher likelihood than the RBF field
with a smaller adjustment index.

The raw-return joint fit provides a corroborating channel comparison. Its
pooled barycentric and news-link coefficients are estimated jointly with a
common intercept, while the factor-conditioned estimates in
\Cref{tab:two-field-pooled} are the primary specification.

\begin{table}[H]
  \centering
  \caption{Raw-return, common-intercept joint barycentric and news-link model.
  The 95\% intervals use 2,000 joint-date stationary-bootstrap refits.}
  \label{tab:two-field-raw-appendix}
  \begin{tabular}{lrrr}
    \toprule
    Channel & $\hat\rho$ & 95\% interval & $\hat\lambda$ \\
    \midrule
    Barycentric ($W^\flat$) &
    \ppnum[3]{post-pooled-2023-2026-two-field-rho-b} &
    $[\ppnum[3]{post-pooled-2023-2026-two-field-rho-b-ci-lower},
    \ppnum[3]{post-pooled-2023-2026-two-field-rho-b-ci-upper}]$ &
    \ppnum[2]{post-pooled-2023-2026-two-field-lambda-b} \\
    News-link ($W^{\mathrm{news}}$) &
    \ppnum[3]{post-pooled-2023-2026-two-field-rho-n} &
    $[\ppnum[3]{post-pooled-2023-2026-two-field-rho-n-ci-lower},
    \ppnum[3]{post-pooled-2023-2026-two-field-rho-n-ci-upper}]$ &
    \ppnum[2]{post-pooled-2023-2026-two-field-lambda-n} \\
    Total &
    \ppnum[3]{post-pooled-2023-2026-two-field-rho-total} &
    $[\ppnum[3]{post-pooled-2023-2026-two-field-rho-total-ci-lower},
    \ppnum[3]{post-pooled-2023-2026-two-field-rho-total-ci-upper}]$ & \\
    \bottomrule
  \end{tabular}
\end{table}

\begin{figure}[htbp]
  \centering
  \ppfigure{1}{two_field_region}
  \caption{Joint two-field likelihood and bootstrap geometry in the raw-return,
    common-intercept specification of \Cref{tab:two-field-raw-appendix}. The horizontal
    axis is barycentric-channel feedback $\rho^\flat$, and the vertical axis is
    news-link feedback $\rho^{\mathrm{news}}$. The left panel places the joint estimate
    inside the admissible triangle, with both single-field fits on the axes and
    the stability boundary $\rho^\flat+\rho^{\mathrm{news}}=1$ marked. The
    right panel is the same
    profiled surface zoomed to the bootstrap cloud, with a descriptive
    two-parameter quasi-log-likelihood contour. The cloud's diagonal
    tilt records
    the trade-off between channels, while its dispersion in both directions
  shows the sampling variation in their allocation.}
  \label{fig:two-field-region}
  \figalttext{Two contour panels plot barycentric-channel feedback horizontally
    and news-link feedback vertically. The left panel shows an interior joint
    estimate within the triangular stability region and single-field estimates
    on the axes. The right panel magnifies an elongated downward-sloping
    bootstrap cloud around the joint estimate, indicating a trade-off in
  allocating feedback between channels.}
\end{figure}

\Cref{fig:two-field-region} shows the profiled joint surface behind
\Cref{tab:two-field-raw-appendix}.
The bootstrap cloud tilts along the diagonal, reflecting the trade-off between
two channels whose peer-return signals correlate at
\ppnum[3]{overlap-induced-pearson}, yet it varies in both directions rather
than collapsing onto a one-dimensional ridge.

\begin{table}[H]
  \centering
  \caption{Annual unadjusted model-implied adjustment under the frozen
    barycentric interaction field. The 2026 period ends on
  \ppvalue{post-evaluation-date-end}.}
  \label{tab:qmle-annual}
  \begin{tabular}{lrrr}
    \toprule
    Evaluation year & $\hat\lambda$ & 95\% interval & Implied $\hat\rho$ \\
    \midrule
    2023 & \ppnum[2]{post-eval2023-w-flat-lambda} &
    $[\ppnum[2]{post-eval2023-w-flat-lambda-ci-lower},
    \ppnum[2]{post-eval2023-w-flat-lambda-ci-upper}]$ &
    \ppnum[3]{post-eval2023-w-flat-rho} \\
    2024 & \ppnum[2]{post-eval2024-w-flat-lambda} &
    $[\ppnum[2]{post-eval2024-w-flat-lambda-ci-lower},
    \ppnum[2]{post-eval2024-w-flat-lambda-ci-upper}]$ &
    \ppnum[3]{post-eval2024-w-flat-rho} \\
    2025 & \ppnum[2]{post-eval2025-w-flat-lambda} &
    $[\ppnum[2]{post-eval2025-w-flat-lambda-ci-lower},
    \ppnum[2]{post-eval2025-w-flat-lambda-ci-upper}]$ &
    \ppnum[3]{post-eval2025-w-flat-rho} \\
    2026 & \ppnum[2]{post-eval2026-w-flat-lambda} &
    $[\ppnum[2]{post-eval2026-w-flat-lambda-ci-lower},
    \ppnum[2]{post-eval2026-w-flat-lambda-ci-upper}]$ &
    \ppnum[3]{post-eval2026-w-flat-rho} \\
    \bottomrule
  \end{tabular}
\end{table}

With the 2018--2022 field held fixed, re-estimating only the common coefficient
and nuisance parameters within each calendar year yields adjustment indices
between \ppnum[2]{post-eval2024-w-flat-lambda} and
\ppnum[2]{post-eval2026-w-flat-lambda}, and all four intervals lie above zero.
Interval widths vary, especially in 2025, and the 2026 estimate uses
\ppvalue{post-eval2026-w-flat-trading-days} dates.
